\section{Introduction}

\subsection{Document Purpose}
This Systems Engineering Plan (SEP) defines the technical engineering approach, processes, tools, and reviews for the development of the Meridian Next-Generation Autonomous Reconnaissance System (M-ARS). This document establishes the authoritative framework for coordinating the multidisciplinary engineering teams across Nexa Defense Systems, Vertex Aerospace, and Synapse Technologies. It ensures compliance with rigorous aerospace development lifecycles and provides a roadmap from initial requirements validation to system qualification.

\subsection{System Scope}
The M-ARS is a long-range tactical surveillance platform deployed for unmanned perimeter protection, search-and-rescue support, and real-time multi-spectral target tracking. The system integrates advanced propulsion, dynamic obstacle avoidance, satellite-backed mesh telecommunication networks, and high-performance edge compute platforms capable of processing onboard computer vision pipelines.

\begin{table}[htbp]
\centering
\caption{M-ARS Subsystem Overview and Interfaces}
\label{tab:system_chars}
\small
\begin{tabularx}{\textwidth}{l p{5cm} X}
\toprule
\textbf{Subsystem} & \textbf{Component Specifications} & \textbf{Key Integration Interface} \\
\midrule
Airframe \& Propulsion & Hybrid electric quad-rotor with vertical takeoff and carbon fiber body & Nexa Core Flight Controller \\
Sensor Payload & Multi-spectral thermal, lidar, and high-definition electro-optical gimbal & Synapse Image Processing Unit \\
Communications & Secured tactical satellite uplink and peer-to-peer RF mesh & Vertex Unified Encryption Module \\
Onboard Compute & Dual-redundant low-latency tensor processing unit & Meridian Data Integration Bus \\
Power Subsystem & 48V Solid-state high-density lithium polymer battery cells & Nimbus Power Management System \\
\bottomrule
\end{tabularx}
\end{table}

\subsection{Program Organization}
The M-ARS engineering team is structured to support concurrent engineering practices. Nexa Defense Systems leads the overall systems integration and flight control software, while Vertex Aerospace manages airframe manufacturing and aeromechanics. Synapse Technologies delivers the advanced sensor payload and onboard machine learning models.

% --- SECTION 2 ---
